\section{Model symulacji pożaru}

\subsection{Przegląd architektury}

Symulator pożaru lasu modeluje fizyczne procesy związane z pożarem, w tym:
\begin{itemize}
    \item Rozprzestrzenianie się ognia między sektorami
    \item Zmiany poziomu ognia w sektorze
    \item Proces spalania (burn level)
    \item Proces gaszenia przez brygady pożarnicze
    \item Wpływ warunków środowiskowych (wiatr, wilgotność, temperatura)
\end{itemize}

\textbf{Ważne:} Symulator działa w module \textbf{FFSim} i jest oddzielony od systemu decyzyjnego. Wszystkie decyzje podejmowane są przez system (FFBackend, FFSup), a symulator jedynie wykonuje komendy i symuluje fizykę.



System \textit{ForestFire} został zaprojektowany jako aplikacja rozproszona. 
Składa się z następujących komponentów:

\begin{table}[h!]
\centering
\caption{Komponenty systemu ForestFire}
\label{tab:forestfire_components}
\begin{tabular}{|l|p{7cm}|l|}
\hline
\textbf{Komponent} & \textbf{Funkcja} & \textbf{Technologia} \\ \hline
FFSup &
Główny moduł analityczny i decyzyjny systemu monitorującego &
Java, Python, Flask \\ \hline
FFVis &
Interfejs wizualizacyjny (panel webowy), służący zarówno dla FFSup, jak i FFConf &
JavaScript, Electron \\ \hline
FFConf &
Moduł konfiguracji środowiska i sektorów (edycja map, punktów, stacji pomiarowych itd.) &
Python, FastAPI \\ \hline
FFSim &
Symulator zdarzeń w monitorowanym obszarze oraz transmisji danych; działa niezależnie, możliwe jest uruchomienie wielu symulatorów &
Python, Flask \\ \hline
Baza danych &
Przechowywanie danych geoprzestrzennych oraz definicji monitorowanych obszarów i sektorów &
MongoDB \\ \hline
\end{tabular}
\end{table}

Komunikacja pomiędzy modułami odbywa się z wykorzystaniem protokołów 
\textit{WebSocket}/RPC oraz brokera wiadomości \textit{RabbitMQ}. 
Zarys architektury systemu przedstawiono na Rys.~1.

\subsection{Architektura fizyczna}

System uruchamiany jest w kontenerach \textit{Docker}. 
Odpowiednio przygotowany skrypt umożliwia uruchomienie całego systemu 
przy pomocy jednego polecenia:

\begin{verbatim}
docker compose -f docker-compose.yml up --build
\end{verbatim}

Każdy moduł może być wdrożony niezależnie, przy pomocy polecenia:

\begin{verbatim}
docker compose -f docker-compose.yml up --build "<nazwa_modułu>"
\end{verbatim}

Gdzie \texttt{<nazwa\_modułu>} to jedna z następujących wartości:

\begin{itemize}
    \item \texttt{rabbitmq-service}
    \item \texttt{forest-simulator-service}
    \item \texttt{mongo}
    \item \texttt{fire-backend-service}
    \item \texttt{fire-configuration-service}
    \item \texttt{fire-visualization-service}
\end{itemize}

W celu poprawnego działania modułu 
\texttt{fire-visualization-service} konieczne jest wcześniejsze
posiadanie skonfigurowanego klucza \textit{Google API} oraz zainstalowanie
zależności środowiska \textit{Electron} za pomocą polecenia:

\begin{verbatim}
npm install
\end{verbatim}

Polecenie to należy wywołać w głównym k


\subsection{Przepływ danych}
\section*{Symulator (FFSim)}
Symulator (FFSim) generuje dane o sytuacji w sektorach i warunkach pogodowych. Wybiera w sposób losowy sektor, w którym rozpocznie się pożar, a następnie rozprzestrzenia go po pobliskich sektorach, biorąc pod uwagę czynniki środowiskowe, takie jak prędkość i kierunek wiatru.  

Symuluje działania strażaków i patroli leśników po odebraniu poleceń z backendu (FFSup).  

Generuje dane z czujników i kamer, a także stan otoczenia, w którym znajdują się patrolujący leśnicy.  

Oblicza rekomendacje rozkazów dla straży pożarnej za pomocą algorytmu Monte Carlo Tree Search, a następnie przesyła je do backendu (FFSup). Przy odpowiednim ustawieniu na frontendzie (FFVis), może automatycznie stosować rekomendacje, bez oczekiwania na zgodę z backendu (FFSup).  

\section*{Backend (FFSup)}
Backend (FFSup) przekazuje dane do wyświetlenia we frontendzie (FFVis), a także przekazuje polecenia dla strażaków i patroli leśników do symulatora (FFSim).  

\section*{Frontend (FFVis)}
Frontend (FFVis) pobiera dane z backendu (FFSup) i wizualizuje je na mapie. Drugą rolą frontend jest pobieranie danych z konfiguratora (FFConf) i wizualizacja procesów edycyjnych monitorowanych obszarów.  

\section*{Administrator (FFConf)}
Administrator (FFConf) – tylko ten moduł może definiować i modyfikować konfigurację map i punktów infrastruktury oraz sektorów leśnych.  

\section*{Maps (Sectors)}
Maps (Sectors) to główna baza danych zawierająca definicje obszarów monitorowanych (mapa podkładowa, sektory, inne). Przechowywane w niej są pliki konfiguracyjne mapy o strukturze opisanej w formacie JSON.  


\subsection{Sektor lasu}

\subsubsection{Stany sektora}




Sektor może znajdować się w jednym z czterech stanów:
\begin{itemize}
    \item \textbf{DORMANT} -- sektor nie płonie (gotowy do zapłonu)
    \item \textbf{BURNING} -- aktywny pożar
    \item \textbf{EXTINGUISHED} -- pożar ugaszony przez brygadę
    \item \textbf{ASH} -- sektor wypalony, paliwo wyczerpane (stan nieodwracalny)
\end{itemize}

Przejścia między stanami:
\begin{itemize}
    \item DORMANT $\rightarrow$ BURNING -- zapłon,
    \item BURNING $\rightarrow$ ASH -- wypalenie paliwa,
    \item BURNING $\rightarrow$ EXTINGUISHED -- skuteczne gaszenie,
    \item EXTINGUISHED $\rightarrow$ ASH -- stan terminalny w przepływie silnika.
\end{itemize}

\subsubsection{Parametry sektora}

Silnik operuje na znormalizowanych parametrach z zakresu $[0,1]$:
\begin{itemize}
    \item \textbf{fireLevel} $\in [0,1]$: intensywność ognia
    \item \textbf{burnLevel} $\in [0,1]$: skumulowany poziom spalenia
    \item \textbf{extinguishLevel} $\in [0,1]$: postęp gaszenia
    \item \textbf{moisture} $\in [0,1]$: wilgotność sektora (obniża prawdopodobieństwo zapłonu)
    \item \textbf{fuel} $\in [0,1]$: dostępne paliwo (zużywane podczas palenia)
    \item \textbf{temperature}: temperatura lokalna sektora (°C)
    \item \textbf{sectorType}: typ sektora (zob. sekcja \ref{sec:alpha})
\end{itemize}

Pozostałe wielkości środowiskowe (wilgotność powietrza, stężenie CO\textsubscript{2}, PM2.5, prędkość i kierunek wiatru) nie są przechowywane na sektorze jako stan symulacji, lecz wyznaczane na bieżąco w warstwie sensorów przy generowaniu telemetrii (zob. sekcja~\ref{sec:sensors-fire}).

\subsection{Aktualizacja poziomu ognia}

Poziom ognia w płonącym sektorze rośnie w każdym ticku o stały przyrost, modyfikowany typem sektora:

\begin{equation}
fireLevel_{t+1} = \mathrm{clamp}\!\left(fireLevel_t + spreadRate \cdot \alpha,\ 0,\ 1\right)
\end{equation}

gdzie:
\begin{itemize}
    \item $spreadRate$ -- bazowy przyrost ognia na tick (domyślnie 0,04)
    \item $\alpha$ -- współczynnik palności typu sektora (zob. sekcja \ref{sec:alpha})
\end{itemize}

Efekty wiatru i temperatury nie wchodzą do tej formuły -- są uwzględniane wyłącznie na etapie prawdopodobieństwa zapłonu sąsiada (sekcja~\ref{sec:ignition}). Wartość 0,04 rozkłada eskalację równomiernie na stopnie klasyfikacji 1--4 (zob. sekcja~\ref{sec:fire-threat-enums}).

\subsection{Aktualizacja poziomu spalenia}

Poziom spalenia (burn level) akumuluje się proporcjonalnie do bieżącej intensywności ognia:

\begin{equation}
burnLevel_{t+1} = \mathrm{clamp}\!\left(burnLevel_t + 0{,}1 \cdot fireLevel_t,\ 0,\ 1\right)
\end{equation}

Niezależnie od poziomu spalenia, o przejściu sektora w stan terminalny decyduje wyczerpanie paliwa: gdy $fuel \leq \varepsilon$, sektor przechodzi w stan \textbf{ASH}.

\subsection{Poziom gaszenia}

Poziom gaszenia nie jest funkcją chwilową, lecz akumuluje się w czasie. Każda brygada obecna w sektorze i będąca w stanie gaszenia podnosi go o stały przyrost na tick:

\begin{equation}
extinguishLevel_{t+1} = extinguishLevel_t + n_{brigades} \cdot extinguishRate
\end{equation}

gdzie $n_{brigades}$ to liczba gaszących brygad, a $extinguishRate$ to przyrost na brygadę (domyślnie 0,05). Gdy $extinguishLevel \geq 1$, pożar uznaje się za ugaszony i sektor przechodzi w stan \textbf{EXTINGUISHED}.

\subsection{Zużycie paliwa}

Paliwo jest zużywane w tempie zależnym od klasyfikacji pożaru (Tabela~1 z artykułu Klimka): im intensywniejszy pożar, tym szybciej wypala paliwo.

\begin{equation}
fuel_{t+1} = \max\!\left(fuel_t - fuelRate \cdot m_{level},\ 0\right)
\end{equation}

gdzie $fuelRate$ to bazowe tempo (domyślnie 0,02), a $m_{level}$ to mnożnik zależny od stopnia pożaru: 0,5 (stopień 1), 1,0 (stopień 2), 2,0 (stopień 3) oraz 4,0 (stopień 4).

\subsection{Wpływ ognia na odczyty sensorów}
\label{sec:sensors-fire}

Silnik nie ewoluuje parametrów środowiskowych jako osobnego stanu sektora. Wpływ pożaru na otoczenie ujawnia się dopiero w warstwie sensorów (faza generacji telemetrii): odczyt sensora stojącego na płonącym sektorze rośnie proporcjonalnie do poziomu ognia $\ell \in [0,1]$. Każdy odczyt to wartość bazowa plus szum (deterministyczny, generowany przez wspólny RNG) plus dorzut od pożaru.

\begin{table}[h!]
\centering
\caption{Wpływ pożaru na odczyty sensorów (przy $\ell = 1{,}0$)}
\begin{tabular}{|l|l|l|}
\hline
\textbf{Sensor} & \textbf{Baseline} & \textbf{Dorzut od pożaru} \\ \hline
Temperatura          & $T_{global}$ & $+45$ °C \\ \hline
Wilgotność powietrza & 0,5          & $-0{,}35$ \\ \hline
Wilgotność ściółki   & 0,3          & $-0{,}25$ \\ \hline
CO\textsubscript{2}  & 400 ppm      & $+3000$ ppm \\ \hline
PM2.5                & 35 µg/m³     & $+450$ µg/m³ \\ \hline
Kamera (dym)         & 5\% szumu    & prawdop. detekcji $\approx 0{,}05 + 0{,}9\,\ell$ \\ \hline
\end{tabular}
\end{table}

Wartości dorzutu dobrano tak, aby płonący sektor wyraźnie przebijał progi wykrywania w silniku (CO\textsubscript{2} $> 800$ ppm, PM2.5 $> 100$ µg/m³, temperatura $> 45$ °C), co stanowi podstawę detekcji pożaru przez sieć czujników (sekcja~\ref{sec:detection}).

\section{Rozprzestrzenianie się ognia}

\subsection{Model rozprzestrzeniania}

Ogień rozprzestrzenia się z sektora aktywnego do sąsiednich sektorów z prawdopodobieństwem zależnym od:
\begin{itemize}
    \item Kierunku i prędkości wiatru
    \item Typu sektora docelowego
    \item Kierunku geograficznego względem wiatru
\end{itemize}

\subsection{Prawdopodobieństwo zapłonu}
\label{sec:ignition}

Dla każdej pary (płonący sąsiad, sektor docelowy) silnik wyznacza prawdopodobieństwo zapłonu jako iloczyn czynników, przeskalowany globalnym współczynnikiem rozprzestrzeniania:

\begin{equation}
p_{ign} = \mathrm{clamp}\!\Big(0,\ 1,\ \ s \cdot f \cdot (1 - m) \cdot w_{\theta} \cdot \ell_{sąsiad} \cdot \max(0,\ 1 + \beta (T - T_{ref}))\Big)
\end{equation}

gdzie:
\begin{itemize}
    \item $s$ -- globalny współczynnik rozprzestrzeniania (domyślnie 0,2); bez niego ogień przy $m=0$ zajmowałby każdego sąsiada deterministycznie,
    \item $f$ -- paliwo sektora docelowego, $m$ -- jego wilgotność,
    \item $w_{\theta}$ -- składowa wiatru (zob. niżej),
    \item $\ell_{sąsiad}$ -- poziom ognia płonącego sąsiada,
    \item $T$ -- temperatura globalna, $T_{ref}$ -- temperatura odniesienia (20 °C), $\beta$ -- współczynnik temperaturowy (0,02).
\end{itemize}

O faktycznym zapłonie decyduje losowanie: sektor zajmuje się ogniem, gdy $\xi < p_{ign}$, $\xi \sim U(0,1)$ ze wspólnego RNG. Sektor świeżo zapalony otrzymuje początkowy $fireLevel = 0{,}1$.

\subsubsection*{Składowa wiatru $w_{\theta}$}

Wiatr nie jest dyskretnym współczynnikiem kierunkowym, lecz ciągłą funkcją kąta między kierunkiem wiatru a kierunkiem rozprzestrzeniania:

\begin{equation}
w_{\theta} = 1 + \alpha_{w} \cdot \|w\| \cdot \cos\theta
\end{equation}

gdzie $\|w\|$ to prędkość wiatru (km/h, bez normalizacji), $\theta$ to kąt między wiatrem a kierunkiem do sektora docelowego, a $\alpha_{w}$ to współczynnik wpływu wiatru (domyślnie 0,01). Wiatr wiejący w stronę sektora ($\cos\theta > 0$) zwiększa prawdopodobieństwo, przeciwny ($\cos\theta < 0$) je zmniejsza.

\subsection{Współczynnik typu sektora $\alpha$}
\label{sec:alpha}

Współczynnik palności $\alpha$ zależy od typu roślinności sektora i wchodzi do przyrostu poziomu ognia (sekcja~\ref{sec:fire-threat-enums}):
\begin{itemize}
    \item \textbf{FIELD:} $\alpha = 1.5$
    \item \textbf{CONIFEROUS:} $\alpha = 1.2$
    \item \textbf{FALLOW:} $\alpha = 1.2$
    \item \textbf{MIXED:} $\alpha = 1.0$
    \item \textbf{UNTRACKED:} $\alpha = 1.0$
    \item \textbf{DECIDUOUS:} $\alpha = 0.8$
    \item \textbf{WATER:} $\alpha = 0.0$ (niepalne)
\end{itemize}

\subsection{Model wiatru}

Wiatr jest globalnym parametrem ewoluującym losowo (faza aktualizacji środowiska). W każdym ticku prędkość i kierunek podlegają niewielkiemu błądzeniu losowemu:

\begin{itemize}
    \item \textbf{Prędkość:} $\|w\|_{t+1} = \mathrm{clamp}(\|w\|_t + \epsilon,\ 0,\ 80)$, $\epsilon \sim U(-1, 1)$ km/h,
    \item \textbf{Kierunek:} $dir_{t+1} = (dir_t + \delta) \bmod 360$, $\delta \sim U(-5, 5)$ stopni.
\end{itemize}

Kierunek wyrażany jest w stopniach (0° = północ, 90° = wschód, 180° = południe, 270° = zachód), a sąsiedztwo jest 4-spójne (północ, wschód, południe, zachód, bez przekątnych).

\section{Baza danych Maps}

Nierelacyjna baza MongoDB przechowuje pliki konfiguracyjne map w formacie JSON. Zawierają one:
\begin{itemize}
    \item Identyfikator i nazwę lasu
    \item Ilość wierszy i kolumn w podziale lasu na sektory
    \item Współrzędne punktów będących wierzchołkami prostokąta reprezentującego las
    \item Listę czujników wraz ze specyfikacją każdego z nich
    \item Listę kamer wraz ze specyfikacją każdego z nich
    \item Listę brygad strażackich wraz z danymi na temat każdej z nich
    \item Listę patroli leśników wraz z danymi na temat każdego z nich
\end{itemize}

Dla każdego sektora zapisane są dodatkowo:
\begin{itemize}
    \item Unikalny identyfikator sektora
    \item Numer wiersza i kolumny, w których znajduje się sektor
    \item Typ sektora
    \item Opis początkowych warunków panujących w sektorze, np. temperatura czy wilgotność
    \item Granice geograficzne sektora, wyrażone w postaci tablicy zawierającej 4 wierzchołki sektora w formacie \([longitude, latitude]\)
\end{itemize}

\section{Format zapisu w bazie danych}

Przykładowy zapis JSON:

\begin{verbatim}
{
   "forestId": "string",
   "forestName": "string",
   "rows": "number",
   "columns": "number",
   "location": [
      {"longitude": "number", "latitude": "number"},
      {"longitude": "number", "latitude": "number"},
      {"longitude": "number", "latitude": "number"},
      {"longitude": "number", "latitude": "number"}
   ],
   "sectors": [
      {
         "sectorId": "number",
         "row": "number",
         "column": "number",
         "sectorType": "enum(SECTOR_TYPE)",
         "initialState": {
            "temperature": "number",
            "windSpeed": "number",
            "windDirection": "enum(GEOGRAPHIC_DIRECTION)",
            "airHumidity": "number",
            "plantLitterMoisture": "number",
            "co2Concentration": "number",
            "pm2_5Concentration": "number"
         },
         "contours": ["number", "number"]
      }
   ],
   "sensors": [
      {
         "sensorId": "number",
         "sensorType": "enum(SENSOR_TYPE)",
         "location": {"longitude": "number", "latitude": "number"},
         "timestamp": "datetime"
      }
   ],
   "cameras": [
      {
         "cameraId": "number",
         "location": {"longitude": "number", "latitude": "number"},
         "range": "number",
         "timestamp": "datetime"
      }
   ],
   "fireBrigades": [
      {
         "fireBrigadeId": "number",
         "timestamp": "datetime",
         "state": "enum(FIREBRIGADE_STATE)",
         "baseLocation": {"longitude": "number", "latitude": "number"},
         "currentLocation": {"longitude": "number", "latitude": "number"}
      }
   ],
   "foresterPatrols": [
      {
         "foresterPatrolId": "number",
         "timestamp": "datetime",
         "state": "enum(FORESTERPATROL_STATE)",
         "baseLocation": {"longitude": "number", "latitude": "number"},
         "currentLocation": {"longitude": "number", "latitude": "number"}
      }
   ]
}
\end{verbatim}

\subsection{Opis pól}

\begin{itemize}
    \item \textbf{forestId} – identyfikator konfiguracji, czyli nazwa lasu.
    \item \textbf{forestName} – nazwa opisowa lasu.
    \item \textbf{rows, columns} – ilość wierszy i kolumn, na które zostanie podzielony obszar lasu.
    \item \textbf{location} – współrzędne geograficzne czterech wierzchołków obszaru lasu.
    \item \textbf{sectors} – lista konfiguracji każdego sektora:
    \begin{itemize}
        \item \textbf{sectorId} – unikalny identyfikator sektora.
        \item \textbf{row, column} – położenie sektora w obrębie lasu, wartości całkowite od 1 do odpowiednio rows i columns.
        \item \textbf{sectorType} – typ sektora; wartość wyliczeniowa SECTOR\_TYPE.
        \item \textbf{initialState} – początkowe warunki w sektorze:
        \begin{itemize}
            \item \textbf{temperature} – temperatura [°C]
            \item \textbf{windSpeed} – prędkość wiatru [m/s]
            \item \textbf{windDirection} – kierunek wiatru, wartość wyliczeniowa GEOGRAPHIC\_DIRECTION
            \item \textbf{airHumidity} – wilgotność powietrza [%]
            \item \textbf{plantLitterMoisture} – wilgotność ściółki [%]
            \item \textbf{co2Concentration} – stężenie CO\(_2\) [ppm]
            \item \textbf{pm2\_5Concentration} – stężenie PM2,5 [µg/m\(^3\)]
        \end{itemize}
        \item \textbf{contours} – granice geograficzne sektora, tablica 4 wierzchołków \([longitude, latitude]\)
    \end{itemize}
    \item \textbf{sensors} – lista czujników:
    \begin{itemize}
        \item \textbf{sensorId} – unikalny identyfikator
        \item \textbf{sensorType} – typ czujnika, wartość wyliczeniowa SENSOR\_TYPE
        \item \textbf{location} – współrzędne geograficzne
        \item \textbf{timestamp} – stempel czasowy
    \end{itemize}
    \item \textbf{cameras} – lista kamer:
    \begin{itemize}
        \item \textbf{cameraId} – unikalny identyfikator
        \item \textbf{location} – współrzędne geograficzne
        \item \textbf{range} – zasięg kamery
        \item \textbf{timestamp} – stempel czasowy
    \end{itemize}
    \item \textbf{fireBrigades} – lista brygad strażackich:
    \begin{itemize}
        \item \textbf{fireBrigadeId} – unikalny identyfikator
        \item \textbf{timestamp} – stempel czasowy
        \item \textbf{state} – aktualna akcja, wartość wyliczeniowa FIREBRIGADE\_STATE
        \item \textbf{baseLocation} – współrzędne bazy strażackiej
        \item \textbf{currentLocation} – aktualne współrzędne brygady
    \end{itemize}
    \item \textbf{foresterPatrols} – lista patroli leśników:
    \begin{itemize}
        \item \textbf{foresterPatrolId} – unikalny identyfikator
        \item \textbf{timestamp} – stempel czasowy
        \item \textbf{state} – aktualna akcja, wartość wyliczeniowa FORESTERPATROL\_STATE
        \item \textbf{baseLocation} – współrzędne bazy grupy leśników
        \item \textbf{currentLocation} – aktualne współrzędne grupy
    \end{itemize}
\end{itemize}

\subsection{API backendu}
Istnieje jeden endpoint HTTP - /run-simulation?interval=<interval>, który pozwala utworzyć subskrypcję na Server Sent Events (SSE) i otrzymywać wiadomości, które wysyłane są z zadanym interwałem czasowym i zawierają aktualny stan symulacji. 


\subsection{Główna komunikacja backend-symulator (FFSup-FFSim) }
Symulator (FFSim) przesyła dane w formacie JSON za pośrednictwem systemu kolejkowego RabbitMQ. Wiadomości te są odbierane przez backend (FFSup) i zawierają następujące informacje:
\begin{itemize}
    \item stan ekip gaśniczych,
    \item stan patroli leśników,
    \item dane z czujników środowiskowych,
    \item dane z kamer monitorujących obszar lasu.
\end{itemize}

\subsubsection{ Diagram przedstawiający kolejki i przepływ danych pomiędzy Forest protection i fire processing a Forest environments simulator}

\subsubsection{Format danych z czujników }

\subsection{Dane stanu czujników}

Dane stanu czujników są publikowane do odpowiednich kolejek brokera RabbitMQ w zależności od typu czujnika. Każdy komunikat posiada część wspólną, identyczną dla wszystkich typów czujników.

\subsubsection*{Część wspólna komunikatu czujnika}

\begin{verbatim}
{
  "sensorId": "number",
  "sensorType": "enum(SENSOR_TYPE)",
  "timestamp": "datetime",
  "location": {
    "longitude": "number",
    "latitude": "number"
  },
  "data": {}
}
\end{verbatim}

\subsubsection*{Czujnik temperatury i wilgotności powietrza}

\begin{verbatim}
{
  "sensorId": "number",
  "timestamp": "datetime",
  "sensorType": "SENSOR_TYPE.TEMPERATURE_AND_AIR_HUMIDITY",
  "location": {
    "longitude": "number",
    "latitude": "number"
  },
  "data": {
    "temperature": "number",
    "airHumidity": "number"
  }
}
\end{verbatim}

\subsubsection*{Czujnik prędkości wiatru}

\begin{verbatim}
{
  "sensorId": "number",
  "timestamp": "datetime",
  "sensorType": "SENSOR_TYPE.WIND_SPEED",
  "location": {
    "longitude": "number",
    "latitude": "number"
  },
  "data": {
    "windSpeed": "number"
  }
}
\end{verbatim}

\subsubsection*{Czujnik kierunku wiatru}

\begin{verbatim}
{
  "sensorId": "number",
  "timestamp": "datetime",
  "sensorType": "SENSOR_TYPE.WIND_DIRECTION",
  "location": {
    "longitude": "number",
    "latitude": "number"
  },
  "data": {
    "windDirection": "enum(GEOGRAPHIC_DIRECTION)"
  }
}
\end{verbatim}

\subsubsection*{Czujnik wilgotności ściółki leśnej}

\begin{verbatim}
{
  "sensorId": "number",
  "timestamp": "datetime",
  "sensorType": "SENSOR_TYPE.LITTER_MOISTURE",
  "location": {
    "longitude": "number",
    "latitude": "number"
  },
  "data": {
    "plantLitterMoisture": "number"
  }
}
\end{verbatim}

\subsubsection*{Czujnik stężenia pyłów PM2.5}

\begin{verbatim}
{
  "sensorId": "number",
  "timestamp": "datetime",
  "sensorType": "SENSOR_TYPE.PM2_5",
  "location": {
    "longitude": "number",
    "latitude": "number"
  },
  "data": {
    "pm2_5Concentration": "number"
  }
}
\end{verbatim}

\subsubsection*{Czujnik stężenia dwutlenku węgla CO\textsubscript{2}}

\begin{verbatim}
{
  "sensorId": "number",
  "timestamp": "datetime",
  "sensorType": "SENSOR_TYPE.CO2",
  "location": {
    "longitude": "number",
    "latitude": "number"
  },
  "data": {
    "co2Concentration": "number"
  }
}
\end{verbatim}

\subsubsection{Format danych z kamer wykrywających dym}

\subsubsection*{Dane stanu kamery}

\begin{verbatim}
{
  "cameraId": "number",
  "timestamp": "datetime",
  "location": {
    "longitude": "number",
    "latitude": "number"
  },
  "data": {
    "smokeDetected": "boolean",
    "smokeLevel": "number",
    "smokeLocation": {
      "longitude": "number",
      "latitude": "number"
    }
  }
}
\end{verbatim}

\subsubsection{Komunikacja Fire Brigade Analyser -- Fire Brigade Simulator}

\subsubsection*{Kierunek: Fire Brigade Analyser $\rightarrow$ Fire Brigade Simulator}

\paragraph{Wysłanie brygady strażackiej do pożaru}

\begin{verbatim}
{
  "fireBrigadeId": "number",
  "action": "FIREBRIGADE_ACTION.EXSTINGUISH",
  "fireState": "enum(FIRE_STATE)",
  "timestamp": "datetime",
  "location": {
    "longitude": "number",
    "latitude": "number"
  }
}
\end{verbatim}

\paragraph{Wysłanie brygady strażackiej do bazy}

\begin{verbatim}
{
  "fireBrigadeId": "number",
  "action": "FIREBRIGADE_ACTION.GO_TO_BASE",
  "timestamp": "datetime"
}
\end{verbatim}

Lokalizacja bazy strażackiej jest pomijana ze względu na wiedzę posiadaną przez każdą brygadę strażacką.

\subsubsection*{Kierunek: Fire Brigade Simulator $\rightarrow$ Fire Brigade Analyser}

\paragraph{Informacja o przemieszczaniu się strażaków}

\begin{verbatim}
{
  "fireBrigadeId": "number",
  "state": "FIREBRIGADE_STATE.TRAVELLING",
  "timestamp": "datetime",
  "location": {
    "longitude": "number",
    "latitude": "number"
  }
}
\end{verbatim}

\paragraph{Informacja o gaszeniu pożaru}

\begin{verbatim}
{
  "fireBrigadeId": "number",
  "state": "FIREBRIGADE_STATE.EXTINGUISHING",
  "timestamp": "datetime",
  "location": {
    "longitude": "number",
    "latitude": "number"
  }
}
\end{verbatim}

\paragraph{Informacja o zakończeniu akcji gaśniczej}

\begin{verbatim}
{
  "fireBrigadeId": "number",
  "state": "FIREBRIGADE_STATE.AVAILABLE",
  "fireState": "FIRE_STATE.EXTINGUISHED",
  "timestamp": "datetime",
  "location": {
    "longitude": "number",
    "latitude": "number"
  }
}
\end{verbatim}

\paragraph{Informacja o dostępności brygady}

\begin{verbatim}
{
  "fireBrigadeId": "number",
  "state": "FIREBRIGADE_STATE.AVAILABLE",
  "timestamp": "datetime",
  "location": {
    "longitude": "number",
    "latitude": "number"
  }
}
\end{verbatim}

\subsubsection{Komunikacja Forester Patrol Analyser -- Forester Patrol Simulator}

\subsubsection*{Kierunek: Forester Patrol Analyser $\rightarrow$ Forester Patrol Simulator}

\paragraph{Wysłanie grupy leśników na patrol}

\begin{verbatim}
{
  "foresterPatrolId": "number",
  "action": "FORESTERPATROL_ACTION.PATROL",
  "timestamp": "datetime",
  "location": {
    "longitude": "number",
    "latitude": "number"
  }
}
\end{verbatim}

\paragraph{Wysłanie grupy leśników do bazy}

\begin{verbatim}
{
  "foresterPatrolId": "number",
  "action": "FORESTERPATROL_ACTION.GO_TO_BASE",
  "timestamp": "datetime"
}
\end{verbatim}

Lokalizacja bazy jest pomijana ze względu na wiedzę posiadaną przez każdą grupę leśników.

\subsubsection*{Kierunek: Forester Patrol Simulator $\rightarrow$ Forester Patrol Analyser}

\paragraph{Informacja o przemieszczaniu się grupy leśników}

\begin{verbatim}
{
  "foresterPatrolId": "number",
  "state": "FORESTERPATROL_STATE.TRAVELLING",
  "timestamp": "datetime",
  "location": {
    "longitude": "number",
    "latitude": "number"
  }
}
\end{verbatim}

\paragraph{Informacja o stanie patrolowanego sektora}

\begin{verbatim}
{
  "foresterPatrolId": "number",
  "state": "FORESTERPATROL_STATE.PATROLLING",
  "timestamp": "datetime",
  "location": {
    "longitude": "number",
    "latitude": "number"
  },
  "sectorState": ""
}
\end{verbatim}

\paragraph{Informacja o dostępności grupy leśników}

\begin{verbatim}
{
  "foresterPatrolId": "number",
  "state": "FORESTERPATROL_STATE.AVAILABLE",
  "timestamp": "datetime",
  "location": {
    "longitude": "number",
    "latitude": "number"
  }
}
\end{verbatim}

\subsection{Komunikacja z modułem wizualizacji}

Moduł wizualizacji (FFVis) pełni rolę usługową dla backendu (FFSup) oraz modułu konfiguracyjnego (FFConf). Komunikacja realizowana jest za pomocą dwukierunkowego połączenia pomiędzy serwisem backendowym napisanym w języku Java z wykorzystaniem frameworku Spring Boot a jednym z procesów aplikacji opartej o framework ElectronJS.

\subsubsection{Format danych w komunikacji pomiędzy modułem przetwarzającym a wizualizującym}

Opis pól w poniższej strukturze danych w formacie JSON jest zgodny z opisem pól zaprezentowanym dla konfiguracji mapy w bazie danych (sekcja~3.2), dlatego nie został tutaj powielony.

\begin{verbatim}
{
  "forestName": "string",
  "timestamp": "datetime",
  "sectors": [
    {
      "sectorId": "number",
      "state": {
        "timestamp": "datetime",
        "temperature": "number",
        "windSpeed": "number",
        "windDirection": "enum(GEOGRAPHIC_DIRECTION)",
        "airHumidity": "number",
        "plantLitterMoisture": "number",
        "co2Concentration": "number",
        "pm2_5Concentration": "number",
        "threatLevel": "enum(THREAT_LEVEL)",
        "fireState": "enum(FIRE_STATE)"
      }
    }
  ],
  "fireBrigades": [
    {
      "fireBrigadeId": "number",
      "sectorId": "number",
      "location": {
        "longitude": "number",
        "latitude": "number"
      },
      "state": "enum(FIREBRIGADE_STATE)",
      "action": "enum(FIREBRIGADE_ACTION)"
    }
  ],
  "foresterPatrols": [
    {
      "foresterPatrolId": "number",
      "timestamp": "datetime",
      "state": "enum(FORESTERPATROL_STATE)",
      "baseLocation": {
        "longitude": "number",
        "latitude": "number"
      },
      "currentLocation": {
        "longitude": "number",
        "latitude": "number"
      }
    }
  ]
}
\end{verbatim}

\section{Model danych i logika systemu}

\subsection{Model encji}

Każdy czujnik posiada unikalny identyfikator, położenie geograficzne oraz wielkość fizyczną, którą mierzy. W zależności od rodzaju mierzonej wielkości, pod kluczem \texttt{data} może znajdować się jedna wartość (np. w przypadku czujnika stężenia CO\textsubscript{2} — stężenie CO\textsubscript{2}) lub wiele wartości (np. w przypadku czujnika temperatury i wilgotności powietrza — temperatura oraz wilgotność).

Jednostki wykonawcze, takie jak brygady straży pożarnej oraz patrole leśników, również posiadają swoje unikalne identyfikatory, aktualne położenie oraz stan operacyjny (np. podróżowanie, gaszenie pożaru, patrolowanie).

\subsubsection{Symulator pożarów -- Forest Environments Simulator}

Struktura modułu \textit{Forest Environments Simulator} odpowiada za symulację środowiska leśnego oraz dynamiki rozprzestrzeniania się pożarów, a także za odwzorowanie działań podejmowanych przez jednostki straży pożarnej i patrole leśników.

Forest Environments Simulator został podzielony na dwa podkomponenty:

\begin{itemize}
    \item \textbf{Sensor Data Generator} – odpowiedzialny za generowanie danych symulujących informacje otrzymywane z rzeczywistych czujników. Moduł ten został podzielony na dalsze podkomponenty, z których każdy odpowiada za konkretny typ czujnika.
    \item \textbf{Fire Brigade Simulator} – odpowiedzialny za odbieranie zleceń akcji strażackich z kolejki \textit{Fire Brigades Action Queue} w brokerze RabbitMQ, symulowanie działań brygad strażackich oraz wysyłanie informacji o stanie brygad do kolejki \textit{Fire Brigade State Queue} w brokerze RabbitMQ.
\end{itemize}

\subsubsection{Aplikacja przetwarzająca -- Forest Protection \& Fire Processing}

Moduł \textit{Forest Protection \& Fire Processing} odpowiada za analizę danych środowiskowych oraz koordynację działań jednostek wykonawczych systemu.

\paragraph{Obsługa brygad strażackich}

Struktura modułu \textit{Forest Protection \& Fire Processing} w zakresie obsługi brygad strażackich obejmuje mechanizmy odpowiedzialne za podejmowanie decyzji operacyjnych, generowanie poleceń akcji oraz monitorowanie stanu realizowanych zadań.


\subsection{Algorytmy wykrywania zagrożeń}
\label{sec:detection}

System analizuje dane środowiskowe w sposób cykliczny, w konfigurowalnym interwale czasowym.

Modelowanie rozprzestrzeniania się pożaru opiera się na modelu probabilistycznym: dla każdego płonącego sektora analizowane są sektory sąsiednie, a prawdopodobieństwo zapłonu wyznacza funkcja \texttt{calculate\_ignition\_probability()} uwzględniająca paliwo, wilgotność, wiatr oraz temperaturę (sekcja~\ref{sec:ignition}).

Detekcja pożaru nie polega na bezpośrednim odczycie stanu sektora, lecz na danych z sieci czujników. Symulator publikuje pełny stan jedynie tych sektorów, które zostały faktycznie wykryte:
\begin{itemize}
    \item przez czujnik na tym samym lub sąsiednim sektorze, którego odczyt przebił próg detekcji (CO\textsubscript{2} $> 800$ ppm, PM2.5 $> 100$ µg/m³, temperatura $> 45$ °C lub kamera wykryła dym) -- dym i stężenia gazów rozchodzą się na sektory sąsiednie, więc czujnik nie musi stać dokładnie na płonącym sektorze, albo
    \item przez patrol leśników stojący na płonącym sektorze lub na sektorze z nim sąsiadującym.
\end{itemize}

Sektory płonące, lecz jeszcze niewykryte, są maskowane w danych dla modułu wsparcia (publikowane jako \texttt{DORMANT}). Dzięki temu logika decyzyjna reaguje dopiero na potwierdzone wykrycie, a nie na wszechwiedzę o stanie mapy.

\subsection{Klasyfikacja stanu pożaru oraz poziom zagrożenia sektora}
\label{sec:fire-threat-enums}

Stan każdego sektora opisują dwa niezależne pola: \texttt{fireState} (faktyczny stan pożaru) oraz \texttt{threatLevel} (poziom zagrożenia pożarowego, niezależny od tego, czy sektor już się pali). Oba pola wyznacza symulator i przekazuje w telemetrii stanu sektora, a backend mapuje je na poniższe enumeracje.

\subsubsection*{Enumeracja \texttt{FIRE\_STATE}}

Wartość wyznaczana jest z ciągłej intensywności pożaru \texttt{fireLevel} $\in [0,1]$:

\begin{itemize}
    \item \texttt{NON\_COMBUSTED} -- sektor bez ognia (\texttt{fireLevel} $= 0$),
    \item \texttt{MILD} -- \texttt{fireLevel} $\in (0,\ 0{,}25]$,
    \item \texttt{MODERATE} -- \texttt{fireLevel} $\in (0{,}25,\ 0{,}5]$,
    \item \texttt{FULL} -- \texttt{fireLevel} $\in (0{,}5,\ 0{,}75]$,
    \item \texttt{SEVERE} -- \texttt{fireLevel} $\in (0{,}75,\ 1{,}0]$,
    \item \texttt{COMBUSTED} -- sektor wypalony (paliwo wyczerpane),
    \item \texttt{EXTINGUISHED} -- sektor ugaszony przez brygadę.
\end{itemize}

Wartość \texttt{NO\_DATA} rezerwowana jest dla sytuacji, w której backend nie otrzymał jeszcze poprawnego stanu sektora.

\subsubsection*{Enumeracja \texttt{THREAT\_LEVEL}}

Poziom zagrożenia przyjmuje wartości \texttt{LOW}, \texttt{MEDIUM}, \texttt{HIGH}, \texttt{VERY\_HIGH} oraz \texttt{CRITICAL} (a także \texttt{NO\_DATA} przy braku danych). Jest to metryka monitoringu leśników, opisująca ryzyko pożarowe sektora.

Dla sektora aktywnie płonącego poziom wynika wprost z klasyfikacji \texttt{fireState}: \texttt{MILD} odpowiada poziomowi \texttt{HIGH}, \texttt{MODERATE} oraz \texttt{FULL} poziomowi \texttt{VERY\_HIGH}, a \texttt{SEVERE} poziomowi \texttt{CRITICAL}. Sektor wypalony lub ugaszony otrzymuje poziom \texttt{LOW}.

Dla sektora bez ognia poziom wyznacza odległość BFS od najbliższego aktywnie płonącego sektora. Symulator oblicza wieloźródłowe BFS startując równocześnie od wszystkich płonących sektorów (zasięg maksymalny 3 sektory):

\begin{itemize}
    \item odległość 1 -- bazowy poziom \texttt{VERY\_HIGH},
    \item odległość 2 -- bazowy poziom \texttt{HIGH},
    \item odległość 3 -- bazowy poziom \texttt{MEDIUM},
    \item poza zasięgiem (odległość $\geq 4$ lub brak pożaru) -- \texttt{LOW}.
\end{itemize}

Bazowy poziom jest następnie modyfikowany przez suchość sektora ($\mathit{dryness} = 1 - \mathit{moisture}$): bardzo suchy sektor ($\mathit{dryness} > 0{,}7$) podnosi poziom o jeden stopień, a wilgotny ($\mathit{dryness} < 0{,}3$) obniża o jeden stopień. Wynik jest obcinany do zakresu \texttt{LOW}--\texttt{CRITICAL}. Dzięki temu zagrożenie tworzy czytelny gradient promieniujący wokół frontu pożaru.

\subsection{Moduł reguł monitorowania oraz akcji gaśniczych}

Moduł reguł wykorzystuje algorytm \textit{Monte Carlo Tree Search} (MCTS), który uruchamiany jest asynchronicznie w osobnym wątku, w konfigurowalnym interwale czasowym.

Algorytm MCTS analizuje sklonowany stan mapy lasu i generuje rekomendowane akcje w postaci par \texttt{(unitId, sectorId)}, reprezentujących przypisanie jednostek gaśniczych do sektorów wymagających interwencji.

Wygenerowane rekomendacje są następnie filtrowane pod kątem dostępności jednostek oraz przesyłane do odpowiednich kolejek komunikacyjnych o wysokim priorytecie.

System obsługuje również akcje manualne, odbierane z kolejek \texttt{Forester Patrol Action Queue} oraz \texttt{Fire Brigades Action Queue} przez menedżera agentów.

\section{Wytwarzanie logów i monitorowanie}

\subsection{Kategorie logów}

System generuje logi w kilku kategoriach, odpowiadających poszczególnym warstwom architektury oraz pełnionym przez nie funkcjom.

\subsubsection*{Logi aplikacyjne}

Logi aplikacyjne obejmują wszystkie zdarzenia zachodzące w backendzie (FFSup), w szczególności:
\begin{itemize}
    \item obsługę kolejek komunikatów (subskrypcje, otrzymywanie wiadomości, błędy),
    \item aktualizację stanu brygad gaśniczych oraz patroli leśników,
    \item obsługę danych pochodzących z czujników i kamer,
    \item reakcje systemu na zdarzenia środowiskowe,
    \item wysyłanie poleceń uruchomienia i sterowania symulacją.
\end{itemize}

Logi te umożliwiają pełne śledzenie przebiegu przetwarzania danych oraz diagnostykę błędów logicznych i komunikacyjnych.

\subsubsection*{Logi symulatora}

Logi symulatora (FFSim) dotyczą w szczególności:
\begin{itemize}
    \item tworzenia oraz konfiguracji kolejek i tematów RabbitMQ,
    \item dodawania producentów i konsumentów komunikatów,
    \item wysyłania i odbierania wiadomości,
    \item inicjalizacji oraz przebiegu pożarów,
    \item zmian stanów sektorów oraz ich wygaszania,
    \item generowanych danych środowiskowych i stanu czujników (np. temperatura, wilgotność).
\end{itemize}

\subsubsection*{Logi frontendu}

Logi frontendu (FFVis) mają głównie charakter techniczny i obejmują:
\begin{itemize}
    \item uruchamianie i zamykanie aplikacji (tryb deweloperski i produkcyjny),
    \item błędy wizualizacji danych,
    \item odbiór i parsowanie danych otrzymanych z backendu.
\end{itemize}

\subsubsection*{Logi systemowe}

Logi systemowe obejmują logi generowane przez:
\begin{itemize}
    \item kontenery Docker,
    \item bazę danych,
    \item brokera RabbitMQ,
    \item technologie użyte do budowy systemu (m.in. Spring Boot, Flask).
\end{itemize}

Logi są gromadzone przez mechanizmy platformy Docker, która zarządza ich rotacją oraz cyklicznym usuwaniem po przekroczeniu określonej liczby uruchomień kontenera lub maksymalnego rozmiaru plików logów.

\subsection{Format logów}

System posiada rozbudowany mechanizm logowania. Logi są wypisywane na standardowe wyjście przez wszystkie podsystemy. Ze względu na zastosowanie konteneryzacji, Docker gromadzi logi osobno dla każdego modułu i zarządza nimi zgodnie z domyślnymi ustawieniami platformy.

Szczegółowy format logów oraz ich przykłady zostały przedstawione w załączniku do dokumentacji.

\subsubsection{Logi modułu symulacji}

\paragraph{Logi powiązane z komunikacją RabbitMQ}

\begin{itemize}
    \item All queues and topics are created and bound.
    \item Queue created: \{queue\_name\}
    \item Queue '\{queue\_name\}' bound to topic '\{topic\_name\}'
    \item Received message \{data\} from \{queue\_name\}
    \item Waiting for messages in queue: \{queue\_name\}
    \item Error consuming messages: \{e\}
    \item Retrieved oldest received message: \{oldest\_message\} from queue: \{queue\_name\}
    \item Message received from Fire brigades queue: \{message\}
    \item Channel is None
\end{itemize}

\paragraph{Logi systemu rekomendacyjnego}

\begin{itemize}
    \item MCTS recommended actions: \{recommended\_actions\}
    \item Error in MCTS prediction: \{str(e)\}
    \item MCTS prediction thread started
    \item Starting MCTS search with time limit \{MAX\_SIMULATION\_TIME\} seconds
\end{itemize}

\paragraph{Logi związane z poziomem ognia}

\begin{itemize}
    \item Fire started in sector \{sector\_id\}, column: \{column\}, row: \{row\}
    \item New fire level in sector \{sector\_id\} is \{fire\_level\}
    \item Sector \{sector\_id\} is extinguished
\end{itemize}

\paragraph{Logi działań brygad gaśniczych}

\begin{itemize}
    \item Order received: Brigade \{fire\_brigade\_id\} returning to base
    \item Order received: Brigade \{fire\_brigade\_id\} extinguishing fire at \{location\}
\end{itemize}

\paragraph{Logi patroli leśników}

Logi dotyczące patroli leśników są analogiczne do logów generowanych dla brygad gaśniczych.

\paragraph{Logi stanu czujników i środowiska}

\begin{itemize}
    \item Camera \{camera\_id\} has data: \{data\}
    \item Initial wind: \{speed\} \{direction\}
\end{itemize}

\subsubsection{Logi backendu}

\paragraph{Logi komunikacji RabbitMQ}

\begin{itemize}
    \item Subscribed to sendMessageToQueue \{\}
    \item Failed to serialize message: \{\}
    \item Failed to declare queue: \{\}
    \item Failed to parse RMQ message: \{\}
\end{itemize}

\paragraph{Logi stanu środowiska}

Przykładowy log:
\begin{quote}
Wind speed updated for sector \{\}: \{\} m/s
\end{quote}

Analogiczne logi istnieją dla wszystkich parametrów środowiskowych.

\paragraph{Logi czujników}

Przykładowy log:
\begin{quote}
Sector at location \{\} not found for temperature and humidity sensor
\end{quote}

Analogiczne logi generowane są dla wszystkich typów czujników, np.:
\begin{quote}
Litter moisture sensor event received at location: \{\}
\end{quote}

\paragraph{Logi stanu symulacji}

\begin{itemize}
    \item Starting simulation with interval: \{\} seconds
    \item Initial state created with \{\} sectors
    \item Sending HTTP request to start simulation
\end{itemize}

\paragraph{Logi brygad gaśniczych}

\begin{itemize}
    \item Fire brigade event received -- ID: \{\}, Location: \{\}
    \item Fire brigade state updated -- ID: \{\}
\end{itemize}

\paragraph{Logi patroli leśników}

Logi dotyczące patroli leśników są analogiczne do logów generowanych dla brygad gaśniczych.

\subsubsection{Logi frontendu}

Logi frontendu mają głównie charakter techniczny i obejmują informacje o uruchomieniu i zatrzymaniu aplikacji oraz o odbiorze i przetwarzaniu danych wizualizowanych w interfejsie użytkownika.

\section{Bezpieczeństwo i niezawodność }
System jest niestabilny. Błędy, takie jak "teleportujące się" lub znikające obiekty występują dość często. Powodują one kolejne błędy, jak np. nieprawidłowe rekomendacje. Momentami można odnieść wrażenie, że na backendzie dzieje się coś innego, niż jest widoczne w warstwie wizualnej. Błędy te wynikają z nieprawidłowej obsługi wątków a także kolejki komunikatów. W bardzo rzadkich sytuacjach możliwe jest duplikowanie obiektów. Niektóre sensory nie działają lub podają dane niezgodne z symulacją. 

\section{Algorytm MCTS}

\subsection{Przegląd}

Monte Carlo Tree Search (MCTS) jest wykorzystywany w module \textbf{FFSup} do generowania rekomendacji decyzyjnych. Algorytm buduje drzewo możliwych stanów systemu, symulując akcje agentów i oceniając nagrody.

\subsection{Moduł działania}

MCTS działa w module \textbf{FFSup} i wykorzystuje następujące wejścia:
\begin{itemize}
    \item Stan wszystkich sektorów (fireLevel, burnLevel, fireState)
    \item Pozycje i stany wszystkich brygad pożarniczych
    \item Pozycje i stany wszystkich patroli leśnych
    \item Konfiguracja lasu (sąsiedztwo sektorów, typy sektorów)
\end{itemize}

Wyjścia:
\begin{itemize}
    \item Lista rekomendowanych akcji: $(agentId, sectorId)$
    \item Każda rekomendacja określa, którą brygadę wysłać do którego sektora
\end{itemize}

\subsection{Faza selekcji (Selection)}

Algorytm wybiera ścieżkę od korzenia do liścia używając formuły UCT (Upper Confidence bound applied to Trees):

\begin{equation}
UCT(node, parent) = \frac{Q[node]}{N[node]} + c \cdot \sqrt{\frac{\ln(N[parent])}{N[node]}}
\end{equation}

gdzie:
\begin{itemize}
    \item $Q[node]$ -- całkowita nagroda dla węzła
    \item $N[node]$ -- liczba wizyt węzła
    \item $N[parent]$ -- liczba wizyt rodzica
    \item $c = \sqrt{2}$ -- współczynnik eksploracji (dla UCT1)
\end{itemize}

Węzeł z najwyższą wartością UCT jest wybierany do ekspansji.

\subsection{Faza ekspansji (Expansion)}

Po dotarciu do liścia, węzeł jest rozwijany przez wygenerowanie wszystkich możliwych akcji:
\begin{itemize}
    \item Dla każdej dostępnej brygady pożarniczej: możliwe akcje to przypisanie do sektora z aktywnym pożarem
    \item Patrolami leśnymi MCTS nie steruje -- ich ruchem zarządza autonomicznie symulator (zob. sekcja \ref{sec:patrol-heuristic})
\end{itemize}

\subsection{Faza symulacji (Simulation)}

Z wybranego węzła wykonywana jest losowa symulacja do stanu terminalnego (lub do maksymalnej liczby kroków). W każdym kroku:
\begin{enumerate}
    \item Wybierana jest losowa akcja dla każdego agenta
    \item Stan systemu jest aktualizowany (rozprzestrzenianie ognia, gaszenie)
    \item Sprawdzane jest, czy osiągnięto stan terminalny
\end{enumerate}

\subsection{Faza propagacji wstecznej (Backpropagation)}

Nagroda z symulacji jest propagowana wstecz do korzenia, aktualizując wartości $Q$ i $N$ dla wszystkich węzłów na ścieżce:

\begin{algorithmic}
\For{węzeł w ścieżce (od liścia do korzenia)}
    \State $N[węzeł] \gets N[węzeł] + 1$
    \State $Q[węzeł] \gets Q[węzeł] + reward$
\EndFor
\end{algorithmic}

\subsection{Strategie nagród}

System wykorzystuje trzy strategie nagród, które różnią się wagami przypisanymi do różnych komponentów:

\subsubsection{Komponenty nagród i kar}

\begin{itemize}
    \item \textbf{LOST\_SECTOR\_PENALTY} -- kara za utracone sektory (stan \texttt{ASH})
    \item \textbf{EXTINGUISHED\_FIRES\_REWARD} -- nagroda za ugaszone pożary
    \item \textbf{SPREAD\_PREVENTION\_REWARD} -- nagroda za zapobieganie rozprzestrzenianiu
    \item \textbf{BURNED\_SECTOR\_PENALTY} -- kara proporcjonalna do poziomu spalenia
    \item \textbf{SECTOR\_FIRE\_LEVEL\_PENALTY} -- kara związana z intensywnością ognia
\end{itemize}

\subsubsection{Strategia agresywna}

Wysokie wagi dla:
\begin{itemize}
    \item EXTINGUISHED\_FIRES\_REWARD
    \item SPREAD\_PREVENTION\_REWARD
\end{itemize}

Niskie wagi dla kar. Skupia się na szybkim gaszeniu pożarów.

\subsubsection{Strategia zrównoważona}

Zbalansowane wagi dla wszystkich komponentów. Optymalizuje zarówno gaszenie, jak i zapobieganie rozprzestrzenianiu.

\subsubsection{Strategia konserwatywna}

Wysokie wagi dla:
\begin{itemize}
    \item SPREAD\_PREVENTION\_REWARD
    \item LOST\_SECTOR\_PENALTY
\end{itemize}

Skupia się na ochronie sektorów przed utratą.

\subsection{Obliczanie nagrody}

Całkowita nagroda dla stanu obliczana jest jako:

\begin{equation}
reward = penalties + rewards
\end{equation}

gdzie:

\begin{equation}
penalties = -lostSectors \cdot LOST\_PENALTY - totalFireIntensity \cdot FIRE\_LEVEL\_PENALTY - totalBurnLevel \cdot BURNED\_PENALTY
\end{equation}

\begin{equation}
rewards = extinguishedCount \cdot EXTINGUISHED\_REWARD + firebreakCount \cdot SPREAD\_PREVENTION\_REWARD
\end{equation}

\subsection{Warunki terminalne}

Symulacja MCTS kończy się, gdy:
\begin{itemize}
    \item Wszystkie pożary zostały ugaszone (żaden sektor nie jest w stanie \texttt{BURNING})
    \item Wszystkie zagrożone sektory zostały utracone (stan \texttt{ASH})
    \item Osiągnięto maksymalną liczbę kroków symulacji
\end{itemize}

\subsection{Ograniczenia obliczeniowe}

\begin{itemize}
    \item \textbf{Czas obliczeń:} domyślnie 1 sekunda na wywołanie MCTS
    \item \textbf{Maksymalna liczba kroków symulacji:} 100 kroków (zabezpieczenie przed nieskończoną pętlą)
    \item \textbf{Liczba workerów MCTS:} konfigurowalna (domyślnie 2)
\end{itemize}

\section{Heurystyki decyzyjne}

\subsection{Autonomiczne patrolowanie leśników}
\label{sec:patrol-heuristic}

Patrolami leśników steruje autonomicznie symulator, a nie moduł wsparcia. Wynika to z zależności typu jajko–kura: support bramkuje wiedzę o pożarze do momentu jego wykrycia (sekcja~\ref{sec:detection}), a wykrycie wymaga m.in. obecności patrolu w pobliżu ognia -- gdyby to support kierował patrolami, nie wysłałby ich tam, gdzie jeszcze nie widzi pożaru. Dlatego rozkazy patrolu przychodzące z modułu wsparcia są w symulatorze ignorowane (gdy autonomiczny patrol jest włączony), a leśnicy objeżdżają teren samodzielnie. Moduł wsparcia koncentruje się na brygadach gaśniczych.

Bezczynny leśnik dostaje cel zgodnie z następującą heurystyką:

\begin{enumerate}
    \item \textbf{Filtrowanie:} brane są tylko sektory palne, nie płonące i nie zajęte przez innego leśnika (cel lub aktualnie patrolowany).
    \item \textbf{Priorytet wieku:} najpierw sektory najdawniej obserwowane przez patrol; sektory nigdy nie patrolowane traktowane są jako najstarsze.
    \item \textbf{Rozpraszanie:} przy równym wieku wybierany jest sektor położony jak najdalej od pozostałych patroli, żeby leśnicy rozeszli się po mapie, a nie zbiegli w jeden rejon.
    \item \textbf{Krótki dojazd:} dopiero w ostateczności rozstrzyga bliskość celu do danego leśnika.
\end{enumerate}

Po kilku tickach postoju w sektorze (\texttt{patrol\_dwell}) leśnik rusza dalej, dzięki czemu patrole stale pokrywają teren i wykrywają pożar zanim ten dorośnie do czujnika. Leśnik znajdujący się na sektorze, który właśnie zapłonął, jest automatycznie wycofywany do bazy (\textit{auto-withdraw}).

\subsection{Heurystyka wymaganych brygad}

Liczba brygad wymaganych do skutecznego gaszenia sektora obliczana jest jako:

\begin{equation}
requiredBrigades = \lceil \frac{fireLevel}{3} \rceil
\end{equation}

\section{Optymalizacje}

\subsection{Ograniczenie przestrzeni akcji}

Aby zmniejszyć przestrzeń akcji w MCTS, system wykorzystuje następujące optymalizacje:

\begin{itemize}
    \item \textbf{TOP\_SECTOR\_PERCENTAGE:} Rozważane są tylko sektory z najwyższym poziomem ognia (domyślnie top 10\%)
    \item \textbf{MAX\_COORDINATED\_SECTORS:} Maksymalna liczba sektorów, do których można przypisać brygady w jednym kroku (domyślnie 10)
\end{itemize}

\subsection{Klonowanie stanu}

Aby uniknąć współdzielenia stanu między węzłami MCTS, wszystkie sektory i agenci są klonowane przed użyciem w symulacji. Zapewnia to poprawność predykcji.

\section{Ograniczenia modelu}

\subsection{Założenia upraszczające}

\begin{itemize}
    \item \textbf{Model wiatru:} Uproszczony model losowych zmian, nie uwzględnia rzeczywistych danych meteorologicznych
    \item \textbf{Rozprzestrzenianie ognia:} Prawdopodobieństwo zapalenia zależy tylko od wiatru i typu sektora, nie uwzględnia topografii
    \item \textbf{Gaszenie:} Liniowa zależność między liczbą brygad a poziomem gaszenia
    \item \textbf{Agenty:} Agenty poruszają się z stałą prędkością, nie uwzględniają przeszkód terenowych
\end{itemize}

\subsection{Ograniczenia obliczeniowe}

\begin{itemize}
    \item MCTS działa w czasie rzeczywistym, więc czas obliczeń jest ograniczony
    \item Symulacja MCTS jest uproszczona w porównaniu do pełnej symulacji fizycznej
    \item Nie wszystkie możliwe akcje są rozważane (optymalizacje przestrzeni akcji)
\end{itemize}

\subsection{Charakter symulacyjny}

Wszystkie dane są generowane przez symulator:
\begin{itemize}
    \item Brak realnych czujników środowiskowych
    \item Brak realnych danych meteorologicznych
    \item Wszystkie parametry są modelowane matematycznie
\end{itemize}

\section{Architektura decyzyjna}

\subsection{Separacja logiki decyzyjnej}

Logika decyzyjna jest oddzielona od agentów. Agent reprezentuje runtime i wykonanie, podczas gdy decyzja jest wymienialnym modułem.

\subsection{Strategie decyzyjne}

System wspiera następujące strategie decyzyjne:

\begin{itemize}
    \item \textbf{Deterministyczna:} Reguły i heurystyki
    \item \textbf{MCTS:} Monte Carlo Tree Search (obecnie używane)
    \item \textbf{MCTS:} Monte Carlo Tree Search (używane w module wsparcia)
    \item \textbf{LLM-driven:} Strategia oparta na Large Language Models (zaimplementowana w agentach)
\end{itemize}

\subsection{LLM jako strategia decyzyjna}

System posiada w pełni zaimplementowaną integrację z modelami LLM (domyślnie \texttt{gpt-4o-mini}) poprzez klasę \texttt{LLMAgentBrain}. Umożliwia to agentom podejmowanie autonomicznych decyzji na podstawie analizy kontekstu w języku naturalnym.

\subsubsection{Architektura LLMAgentBrain}

Każdy agent (strażak, patrol) posiada własny moduł \texttt{LLMAgentBrain}, który odpowiada za:
\begin{enumerate}
    \item \textbf{Percepcję:} Zbieranie informacji o otoczeniu (sektory z pożarami), stanie własnym agenta oraz działaniach innych agentów (peer announcements).
    \item \textbf{Kontekst:} Budowanie promptu systemowego i użytkownika. Prompt systemowy definiuje rolę agenta, a prompt użytkownika zawiera aktualny stan świata sformatowany jako tekst.
    \item \textbf{Decyzję:} Komunikację z API OpenAI (przez \texttt{LLMClient}) w celu uzyskania decyzji w formacie JSON.
\end{enumerate}

\subsubsection{Przepływ informacji w LLM}

Proces decyzyjny LLM wygląda następująco:
\begin{itemize}
    \item \textbf{Input:} Pozycja agenta, stan paliwa/wody, lista aktywnych pożarów, warunki pogodowe (wiatr), ostatnie komunikaty od innych agentów.
    \item \textbf{Przetwarzanie:} Model językowy analizuje sytuację, biorąc pod uwagę odległość do pożarów, ich priorytet oraz działania innych jednostek, aby uniknąć duplikacji pracy.
    \item \textbf{Output:} Struktura JSON zawierająca typ akcji (np. \texttt{MOVE}, \texttt{EXTINGUISH}) oraz parametry (np. \texttt{sectorId}). Model generuje również \texttt{reasoning} - uzasadnienie decyzji w języku naturalnym.
\end{itemize}

\subsubsection{Klient LLM i Konfiguracja}

Komunikacja odbywa się poprzez asynchronicznego klienta \texttt{LLMClient}. System wymaga poprawnej konfiguracji klucza API (\texttt{OPENAI\_API\_KEY}). W przypadku braku dostępu do API, agent nie podejmuje decyzji autonomicznie (fallback do braku akcji lub sterowania ręcznego).


\section{Przepływ decyzyjny}

\subsection{Cykl decyzyjny}

\begin{enumerate}
    \item \textbf{Inicjalizacja:} System uruchamia symulację (\texttt{SimpleSimulationEngine}).
    \item \textbf{Symulacja:} Silnik aktualizuje stan świata (pożary, wiatr) i agentów.
    \item \textbf{Wsparcie (MCTS):} \texttt{SupportService} w tle analizuje stan lasu i generuje globalne rekomendacje przy użyciu MCTS, które są wysyłane przez RabbitMQ.
    \item \textbf{Autonomia (LLM):} Równolegle, agenci posiadający \texttt{LLMAgentBrain} okresowo (co ok. 10s) analizują otoczenie i podejmują autonomiczne decyzje o ruchu lub gaszeniu.
    \item \textbf{Interakcja:} Użytkownik widzi stan i rekomendacje w \texttt{FFVis}, mogąc wydawać rozkazy ręczne, które nadpisują autonomię agentów.
    \item \textbf{Wykonanie:} Silnik aplikuje decyzje (ruch agentów, gaszenie) i aktualizuje fizykę w kolejnych krokach czasowych.
\end{enumerate}

\subsection{Moduły odpowiedzialne}

\begin{itemize}
    \item \textbf{FFSup:} Generowanie rekomendacji (MCTS, LLM, heurystyki)
    \item \textbf{FFBackend:} Agregacja i publikacja rekomendacji
    \item \textbf{FFVis:} Prezentacja rekomendacji użytkownikowi
    \item \textbf{FFSim:} Wykonanie akcji i aktualizacja stanu fizycznego
\end{itemize}

\section{Infrastruktura Systemowa}

\subsection{Komunikacja (RabbitMQ)}
System wykorzystuje architekturę zdarzeniową opartą o kolejki wiadomości RabbitMQ, co zapewnia luźne powiązanie między modułami symulacji i wsparcia. Kluczowe kanały komunikacji obejmują:

\begin{itemize}
    \item \textbf{support.data.aggregated:} Publikowane przez symulator, zawiera zagregowany stan świata (sektory, agenty) niezbędny dla modułu wsparcia.
    \item \textbf{simulation.recommendations:} Publikowane przez moduł wsparcia, zawiera rekomendacje działań (np. przemieszczenie brygady) wygenerowane przez MCTS.
    \item \textbf{fire\_brigade\_actions / forester\_actions:} Kanały dla konkretnych rozkazów kierowanych do agentów.
\end{itemize}

\subsection{Zarządzanie Stanem}
Moduł \textbf{FFSup} wykorzystuje klasę \texttt{StateManager} do utrzymania spójnego obrazu świata:
\begin{enumerate}
    \item \textbf{Wątkobezpieczeństwo:} Wszystkie operacje odczytu i zapisu są chronione blokadami (locks), co pozwala na równoległe przetwarzanie wiadomości z RabbitMQ i zapytań od workerów MCTS.
    \item \textbf{Minimalne aktualizacje:} System przetwarza i przechowuje tylko niezbędne zmiany stanu, co optymalizuje wydajność przy dużej liczbie sektorów.
    \item \textbf{Snapshoty:} Udostępnia mechanizm pobierania pełnej kopii stanu (snapshot) dla algorytmów decyzyjnych, gwarantując, że MCTS i LLM operują na niezmiennych w czasie obliczeń danych.
\end{enumerate}

\section{Przyszłe rozszerzenia}

\subsection{Planowane ulepszenia algorytmów}

\begin{itemize}
    \item \textbf{FFDI (Fire Weather Index):} Wprowadzenie wskaźnika pogody pożarowej do modelu rozprzestrzeniania
    \item \textbf{Zaawansowane modele wiatru:} Uwzględnienie rzeczywistych danych meteorologicznych
    \item \textbf{Topografia:} Wpływ ukształtowania terenu na rozprzestrzenianie ognia
    \item \textbf{Dynamiczne priorytety:} Adaptacyjne wagi w strategiach nagród
\end{itemize}

\subsection{Rozwój Modelu Językowego}

\begin{itemize}
    \item \textbf{Fine-tuning:} Dostrojenie modelu do specyfiki domeny pożarniczej w celu zmniejszenia liczby tokenów i zwiększenia precyzji.
    \item \textbf{Memory-augmentation:} Wprowadzenie długoterminowej pamięci dla agentów, aby pamiętali historię pożarów w danym sektorze.
\end{itemize}

\section{Podsumowanie}

System wykorzystuje zaawansowane algorytmy (MCTS) i heurystyki do generowania rekomendacji decyzyjnych. Kluczowe cechy:

\begin{itemize}
    \item Separacja symulatora od systemu decyzyjnego
    \item Hybrydowa logika decyzyjna łącząca globalną optymalizację MCTS z lokalną autonomią agentów opartą o LLM
    \item Model fizyczny symulacji z uwzględnieniem wiatru i typu sektora
    \item System nagród i kar z możliwością konfiguracji strategii
    \item Optymalizacje przestrzeni akcji dla wydajności
\end{itemize}

